\cotefile{105}
\scanpage{2}
\unsure{33} -- NODE E33

\margin{\unsure{S E333}\quad $\vert$}

\unsure{$\mathcal Z_a$}

\diagram{\begin{tikzcd}[row sep=small,column sep=large]
\partial \arrow[r,"\unsure{i}"] \arrow[d,phantom,"\big\vert" description] & I \arrow[d,phantom,"\unsure{0\text{-connexe}}" description] \\
\deleted{\unsure{\operatorname{card}}}(\pi_0\partial)\geq 2 & \\
\deleted{\unsure{0\text{-connexe}}} &
\end{tikzcd}}

\[
\left\{\begin{array}{l}
\varnothing\to e\\[2pt]
\partial\times S^{\unsure{n}}(N)\to I\times S^{\unsure{n}}(N)
\end{array}\right.
\]
\inserted{+ Cônes \unsure{gauches} (compl. droits) sur les précédents (donc ici $\{0\}\to\Delta^1$ \unsure{puis} $\{1\}\to\Delta^1$)}

\uline{$\operatorname{Hom}$}$(\unsure{T},X)$

\uline{Lemme} a) Si $X\to e$ a la propriété RLP \unsure{Rel.} pour $\varnothing\to e$ \uline{et} les $\partial S^{\unsure{n}}(N)\to I\times S^{\unsure{n}}(N)$, $N\in\mathcal N$ (partie infinie de $\mathbf N$), alors $X$ est $W_\infty$-asph.

b) Si $X\to Y$ a la RLP p.r. aux cônes gauches (compl. droits) p.r. aux morphismes précédents, alors $X$ est $W_\infty$-propre (resp. $W_\infty$-lisse).

\uline{Cor.} a) Si $X\to Y$ a la propriété RLP p.r. aux \unsure{imm.} de $\unsure{\Delta}$ \deleted{\unsure{imm.}\ \illegible{1}} \inserted{entre catégories \unsure{finies}} \inserted{(compl. \unsure{finies})} \deleted{\illegible{3}}, alors c'est une \unsure{fibration} $W_\infty$-propre (compl. lisse) \inserted{$W_\infty$-fibration \unsure{asphérique}}.

b) Si $f\colon X\to Y$ a la propriété RLP p.r. aux \unsure{imm.} \deleted{\unsure{d'ensembles ordonnés finis}\ \unsure{homol.}\ \illegible{1}} \unsure{de catégories} \inserted{(de cat. finies)} qui sont des $W_\infty$, alors $f$ est $W_\infty$-propre (resp. $W_\infty$-lisse).
S'il a la dite propriété pour les $W_\infty$-cofibrations \inserted{entre cat. finies} \unsure{triviales}, alors \unsure{est} $W_\infty$-parfait.

\mspage{1}

